\documentclass[10pt,a4paper]{amsart}
\usepackage{amsmath,amssymb,amsthm,graphicx}
% For commutative diagrams we use Xy-pic
\usepackage[all]{xy}

% We use 2cell for 2-commutative diagrams.
\xyoption{2cell}
% Theorem environments.
%
\theoremstyle{plain}
\newtheorem{theorem}[subsection]{Theorem}
\newtheorem{proposition}[subsection]{Proposition}
\newtheorem{lemma}[subsection]{Lemma}

\theoremstyle{definition}
\newtheorem{definition}[subsection]{Definition}
\newtheorem{example}[subsection]{Example}
\newtheorem{exercise}[subsection]{Exercise}
\newtheorem{situation}[subsection]{Situation}

\theoremstyle{remark}
\newtheorem{remark}[subsection]{Remark}
\newtheorem{remarks}[subsection]{Remarks}

\numberwithin{equation}{subsection}

% Macros
%
\def\lim{\mathop{\mathrm{lim}}\nolimits}
\def\colim{\mathop{\mathrm{colim}}\nolimits}
\def\Spec{\mathop{\mathrm{Spec}}}
\def\Hom{\mathop{\mathrm{Hom}}\nolimits}
\def\Ext{\mathop{\mathrm{Ext}}\nolimits}
\def\SheafHom{\mathop{\mathcal{H}\!\mathit{om}}\nolimits}
\def\SheafExt{\mathop{\mathcal{E}\!\mathit{xt}}\nolimits}
\def\Sch{\mathit{Sch}}
\def\Mor{\mathop{\mathrm{Mor}}\nolimits}
\def\Ob{\mathop{\mathrm{Ob}}\nolimits}
\def\Sh{\mathop{\mathit{Sh}}\nolimits}
\def\NL{\mathop{N\!L}\nolimits}
\def\CH{\mathop{\mathrm{CH}}\nolimits}
\def\proetale{{pro\text{-}\acute{e}tale}}
\def\etale{{\acute{e}tale}}
\def\QCoh{\mathit{QCoh}}
\def\Ker{\mathop{\mathrm{Ker}}}
\def\Im{\mathop{\mathrm{Im}}}
\def\Coker{\mathop{\mathrm{Coker}}}
\def\Coim{\mathop{\mathrm{Coim}}}

% Boxtimes
%
\DeclareMathSymbol{\boxtimes}{\mathbin}{AMSa}{"02}

%
% Macros for moduli stacks/spaces
%
\def\QCohstack{\mathcal{QC}\!\mathit{oh}}
\def\Cohstack{\mathcal{C}\!\mathit{oh}}
\def\Spacesstack{\mathcal{S}\!\mathit{paces}}
\def\Quotfunctor{\mathrm{Quot}}
\def\Hilbfunctor{\mathrm{Hilb}}
\def\Curvesstack{\mathcal{C}\!\mathit{urves}}
\def\Polarizedstack{\mathcal{P}\!\mathit{olarized}}
\def\Complexesstack{\mathcal{C}\!\mathit{omplexes}}
% \Pic is the operator that assigns to X its picard group, usage \Pic(X)
% \Picardstack_{X/B} denotes the Picard stack of X over B
% \Picardfunctor_{X/B} denotes the Picard functor of X over B
\def\Pic{\mathop{\mathrm{Pic}}\nolimits}
\def\Picardstack{\mathcal{P}\!\mathit{ic}}
\def\Picardfunctor{\mathrm{Pic}}
\def\Deformationcategory{\mathcal{D}\!\mathit{ef}}
\newenvironment{reference}{}{}
\newenvironment{slogan}{}{}
\newenvironment{history}{}{}
\title{Injective maps from a unimodular lattice modulo an orthogonal direct sum onto two dual quotients}
\author{The Stacks Project Authors}
\date{stacks-project:models.tex:local-snapshot-20261012; GNU FDL 1.2 or later}
\begin{document}
\maketitle
\noindent\textit{Editorial note.} The counted unit is the Stacks Project Lemma lemma-orthogonal-direct-sum at source lines 423-438 together with its proof at source lines 440-459, a straight adjacent pairing in the chapter file. The chapter is the upstream file models.tex, which is the file name recorded in the source record; the pool directory carries the longer name stacks-models and is not used as the chapter name anywhere in this card.\par\medskip
\noindent\textit{Editorial note.} This chapter's candidate index records the statement body and the proof body without their closing markers, so the retained spans used here are the full environments: the statement span ends at the line of the closing marker of the lemma environment and the proof span starts at the opening marker of the proof environment. The same one-line accounting was recorded in the earlier card of this chapter, P2-models-recurring-real.\par\medskip
\noindent\textit{Editorial note.} Neither the retained statement nor the retained proof contains a reference to another labelled result of the chapter and neither contains a citation, so this card carries no navigation pointer and no bibliography entry. Two pieces of the author's own notation, the sharp superscript for a dual and the angle brackets for the form, are defined inside the retained statement span itself.\par\medskip
\noindent\textit{Editorial note.} The retained preamble is the chapter preamble's Xy-pic setup at lines 14-18, the theorem declarations and mathematical macros at lines 154-217, and the three no-op wrapper environments. The Xy-pic lines are kept although this particular span draws nothing, because the retained preamble block is taken verbatim and as a whole.\par\medskip
\noindent\textit{Editorial note.} The proof is the author's own complete argument and contains no omission wording: the orthogonal decomposition, the integrality of the projection, the four-term exact sequence and the division by the orthogonal direct sum are all written out.\par\medskip
\noindent\textit{Allowed standard prerequisites (editorial).}\begin{itemize}
\item Integral symmetric bilinear forms on lattices and their duals. For a finite free module over the integers carrying an integral symmetric bilinear form the dual of a submodule is the set of vectors of the rational span pairing integrally with it, written with a sharp superscript. The form is called unimodular when the module equals its own dual, and nondegeneracy on a summand follows from positive definiteness of the form; the author's own notation for the dual and for the form is retained inside the counted span.. The counted lemma is the technical heart of the surrounding block of the chapter: it compares the quotient of a lattice by an orthogonal direct sum with the two torsion quotients attached to a submodule and its orthogonal complement.
\end{itemize}\medskip
\begin{lemma}
\label{lemma-orthogonal-direct-sum}
Let $L$ be a finite free $\mathbf{Z}$-module endowed
with an integral symmetric bilinear positive definite
form $\langle\ ,\ \rangle : L \times L \to \mathbf{Z}$.
Let $A \subset L$ be a submodule with $L/A$ torsion free. Set
$B = \{b \in L \mid \langle a, b\rangle = 0,\ \forall a \in A\}$.
Then we have injective maps
$$
A^\#/A \leftarrow L/(A \oplus B) \rightarrow B^\#/B
$$
whose cokernels are quotients of $L^\#/L$. Here
$A^\# = \{a' \in A \otimes \mathbf{Q} \mid
\langle a, a'\rangle \in \mathbf{Z},\ \forall a \in A\}$
and similarly for $B$ and $L$.
\end{lemma}

\begin{proof}
Observe that
$L \otimes \mathbf{Q} = A \otimes \mathbf{Q} \oplus B \otimes \mathbf{Q}$
because the form is nondegenerate on $A$ (by positivity).
We denote $\pi_B : L \otimes \mathbf{Q} \to B \otimes \mathbf{Q}$
the projection. Observe that $\pi_B(x) \in B^\#$ for $x \in L$
because the form is integral. This gives an exact sequence
$$
0 \to A \to L \xrightarrow{\pi_B} B^\# \to Q \to 0
$$
where $Q$ is the cokernel of $L \to B^\#$. Observe that $Q$
is a quotient of $L^\#/L$ as the map $L^\# \to B^\#$ is surjective
since it is the $\mathbf{Z}$-linear dual to $B \to L$ which is split
as a map of $\mathbf{Z}$-modules.
Dividing by $A \oplus B$ we get a short exact sequence
$$
0 \to L/(A \oplus B) \to B^\#/B \to Q \to 0
$$
This proves the lemma.
\end{proof}

\end{document}
